% AY2023 材料力学 〔I〕（転記）
% 出典: 04_材料力学/original/04_材料力学/2023/2023za.pdf
% 要約: 剛体でできた断熱板に接合された長さ2Lの棒\textcircled{1}, 長さLの棒\textcircled{2}が壁に固定。
%       断熱板の中央(点B)に外力Pを負荷。棒\textcircled{1}:断面積2S,E,α; 棒\textcircled{2}:断面積S,3E,4α。
%       (1) σ1,σ2, (2) λ1,λ2,
%       (3) Pを負荷したまま棒\textcircled{2}のみΔT上昇で各棒が初期長さに戻る → ΔTをPで表す,
%       (4) このときの応力σ'1,σ'2。
%       (温度上昇前のλ1,λ2はLに対し十分小, 棒\textcircled{1}の温度は変化しない)

\section*{〔I〕}

図1のように, 剛体でできた断熱板に接合された長さ $2L$ の棒\textcircled{1}および長さ $L$ の棒\textcircled{2}が
壁に固定されており, 断熱板の中央（点 B）に外力 $P$ を負荷する場合を考える.
棒\textcircled{1}, 棒\textcircled{2}それぞれの断面積を $2S$, $S$, 縦弾性係数を $E$, $3E$, 線膨張係数を
$\alpha$, $4\alpha$ とするとき, 次の問い (1)\,$\sim$\,(4) に答えよ.

\begin{center}
\begin{tikzpicture}[>=Latex]
  % 左壁, 右壁
  \fill[pattern=north east lines] (-0.5,-1.4) rectangle (0,1.4); \draw (0,-1.4) -- (0,1.4);
  \fill[pattern=north east lines] (10,-1.4) rectangle (10.5,1.4); \draw (10,-1.4) -- (10,1.4);
  % 棒\textcircled{1}（A-板, 太, 2S, 半高0.7）
  \draw (0,0.7) -- (6,0.7); \draw (0,-0.7) -- (6,-0.7);
  % 断熱板（B, 縦長方形）
  \fill[gray!20] (6,-1.0) rectangle (6.4,1.0); \draw (6,-1.0) rectangle (6.4,1.0);
  % 棒\textcircled{2}（板-C, 細, S, 半高0.35）
  \draw (6.4,0.35) -- (10,0.35); \draw (6.4,-0.35) -- (10,-0.35);
  % B点 と 外力P（右向き）
  \fill (6.2,0) circle (1.8pt); \node[below] at (6.2,-0.05) {B};
  \draw[->,very thick] (6.4,0.0) -- (7.6,0.0) node[right]{$P$};
  % 節点
  \node[above right] at (0,-0.7) {A}; \node[below left] at (10,-0.35) {C};
  % ラベル
  \node[above] at (3.0,0.7) {棒\textcircled{1}}; \node[above] at (8.2,0.35) {棒\textcircled{2}}; \node[below] at (6.2,-1.0) {断熱板};
  % 寸法 2L, L
  \draw[<->] (0,1.15) -- (6,1.15) node[midway,fill=white]{$2L$};
  \draw[<->] (6.4,1.15) -- (10,1.15) node[midway,fill=white]{$L$};
  % x軸
  \draw[->] (0,-1.05) -- (1.3,-1.05) node[right]{$x$};
\end{tikzpicture}

\vspace{1mm}
図1
\end{center}

棒\textcircled{1}, 棒\textcircled{2}に生じる応力, 変位を求めた後, 次に, 外力 $P$ を負荷したまま棒\textcircled{2}のみ温度を
$\Delta T$ 上昇させると, 各棒は初期長さに戻った. 温度上昇前における棒\textcircled{1}の変位 $\lambda_1$ と
棒\textcircled{2}の変位 $\lambda_2$ は $L$ に対して十分小さく, 棒\textcircled{1}の温度は変化しないものとする.

\begin{enumerate}[label=(\arabic*)]
  \item 棒\textcircled{1}と棒\textcircled{2}に生じる応力 $\sigma_1$ と $\sigma_2$ を求めよ.

  \item 棒\textcircled{1}と棒\textcircled{2}の変位 $\lambda_1$ と $\lambda_2$ を求めよ.

  \item $\Delta T$ を外力 $P$ を用いて答えよ.

  \item 棒\textcircled{1}と棒\textcircled{2}に生じる応力 $\sigma'_1$ と $\sigma'_2$ を求めよ.
\end{enumerate}
